\chapter{Evaluación de Enfoques Éticos}
\label{ch:etica}

\section{Privacidad de datos de clientes}

El microservicio de clientes almacena datos personales identificables: nombre, correo
electrónico, dirección y comuna. Al tratarse de un sistema que atiende tanto canal
presencial como online, estos datos circulan entre varios microservicios (por ejemplo,
\texttt{envíos} necesita la dirección de destino asociada a un pedido).

Frente a esto se aplicó un principio de \textbf{minimización de datos}: cada
microservicio conoce únicamente los campos estrictamente necesarios para su función. El
microservicio de pagos, por ejemplo, no tiene acceso al nombre ni a la dirección del
cliente -- solo al identificador del pedido y al monto a cobrar --, de modo que un
eventual incidente de seguridad en el módulo de pagos no expondría por sí solo datos
personales identificables.

Inicialmente el sistema no implementaba autenticación ni cifrado de credenciales,
decisión de alcance consciente dado el tiempo disponible del proyecto y no un descuido.
Sin embargo, precisamente por tratarse de datos personales de clientes -- direcciones,
historial de compras --, se decidió cerrar esa brecha al menos en el punto donde más
importa: el microservicio de clientes ahora exige un token JWT válido para modificar o
eliminar un perfil, y las contraseñas se almacenan con \texttt{BCryptPasswordEncoder}
(hash con salt, nunca en texto plano). La protección se limitó deliberadamente a las
operaciones de escritura sobre el propio perfil (\texttt{PUT}/\texttt{DELETE}) y no se
extendió a las lecturas ni al resto de los microservicios: \texttt{carrito} y
\texttt{reseñas} consultan clientes vía Feign para sus propias validaciones cruzadas, y
esas llamadas de servicio a servicio no llevan el token de ningún usuario. Proteger
también las lecturas habría exigido propagar un token de servicio a través de Feign --
una solución más completa, pero fuera del alcance de tiempo disponible. Esta es, en sí
misma, una decisión ética explícita: priorizar cerrar la brecha más sensible (que
cualquiera pueda modificar los datos de otro cliente) por sobre una implementación
completa pero más costosa.

\section{Responsabilidad en el despliegue}

BookPoint Chile opera simultáneamente en tres sucursales físicas y un canal web. A
diferencia de un negocio puramente digital, una caída del sistema no solo afecta las
ventas online: si el monolito actual falla, se detiene también la venta presencial en
Concepción, Temuco y La Serena, porque todo depende del mismo proceso.

La arquitectura de microservicios propuesta mitiga parcialmente este riesgo mediante
\textbf{aislamiento de fallas}: si el microservicio de reseñas deja de responder, el
cliente puede seguir explorando el catálogo, agregando productos al carrito y pagando
su compra sin verse afectado. Esto traslada la responsabilidad del despliegue desde "no
puede fallar nunca" (inviable en cualquier sistema real) hacia "una falla puntual no
debe paralizar el negocio completo" -- un estándar más realista y, a la vez, más
exigente en el diseño de cada servicio individual, porque cada uno debe manejar
adecuadamente los casos en que sus dependencias no respondan (ver
sección~\ref{ch:rest}, manejo de errores en comunicación remota).

\section{Impacto en el empleo}

El sistema automatiza tareas hoy realizadas manualmente por el personal de BookPoint:
el registro de ventas en caja (equivalente al microservicio de pedidos), la consulta de
disponibilidad de stock entre sucursales (inventario) y el seguimiento del estado de un
envío (envíos, con su historial de estados).

Es importante encuadrar este punto con honestidad: el sistema no reemplaza al
Asistente de Ventas, al Encargado de Bodega ni al Área de Logística -- automatiza
tareas repetitivas y propensas a error humano (como mantener sincronizado el stock
entre tres sucursales a mano), liberando tiempo para tareas que sí requieren juicio
humano: atención al cliente, resolución de casos excepcionales, gestión de proveedores.
El riesgo ético real no está en automatizar estas tareas puntuales, sino en cómo se
gestione la transición del personal hacia esas nuevas responsabilidades -- una decisión
organizacional de BookPoint Chile que excede el alcance técnico de este informe, pero
que vale la pena nombrar explícitamente como parte de la evaluación ética del proyecto.

\section{Desafíos éticos específicos enfrentados durante el desarrollo}

Más allá de los principios generales discutidos arriba, el desarrollo del proyecto
presentó tensiones concretas entre avanzar rápido y hacerlo bien, y al menos una duda
genuina sobre si una decisión de alcance era aceptable para un sistema real.

\textbf{Errores que exigieron parar y corregir en vez de seguir avanzando.} En varios
momentos del desarrollo surgieron problemas que hubiera sido más rápido ignorar o
parchar superficialmente, pero que se decidió corregir de raíz:

\begin{itemize}
    \item Durante la mayor parte del desarrollo, cualquier recurso no encontrado
    devolvía \texttt{500} en lugar de \texttt{404} en los diez servicios. Era un error
    silencioso -- el sistema "funcionaba" en el sentido de que no se caía --, pero
    devolver el código HTTP correcto es justamente lo que permite que un cliente de la
    API distinga un error suyo de un error del servidor. Se optó por parar y corregirlo
    en los diez servicios a la vez, en lugar de dejarlo para "después".
    \item Un cliente de la API podía enviar un \texttt{id} propio al crear un recurso, y
    ese \texttt{id} terminaba sobrescribiendo silenciosamente un recurso existente en
    vez de crear uno nuevo. Es el tipo de bug que puede pasar desapercibido en las
    demos (todo parece funcionar) pero que en un sistema real significaría corromper
    datos de otro cliente sin ningún aviso.
    \item Al centralizar la documentación Swagger de los diez microservicios a través
    del Gateway, las rutas no llegaban correctamente: hubo que re-enrutar manualmente
    los puertos en el YAML del Gateway y ajustar la configuración de cada
    microservicio para que calzara con lo que el Gateway esperaba. Y al containerizar
    los servicios, Swagger dejó de funcionar porque cada servicio reportaba su propia
    IP interna de Docker como servidor, en vez de una URL que el navegador del cliente
    pudiera alcanzar. Ninguno de los dos problemas era obvio a primera vista, y ambos
    se podrían haber dejado "así no más" con la excusa de que la API seguía
    respondiendo -- se optó por diagnosticar la causa real en vez de silenciar el
    síntoma.
    \item La decisión de alcance de la autenticación JWT (sección~\ref{ch:etica}) es en
    sí misma un ejemplo de esta tensión: la alternativa más rápida habría sido proteger
    todo el microservicio de clientes sin pensar en las consecuencias; se optó por
    invertir tiempo en entender qué llamadas dependían de un acceso sin token
    (\texttt{carrito} y \texttt{reseñas} vía Feign) antes de decidir el alcance
    correcto de la protección.
\end{itemize}

\textbf{Duda genuina sobre decisiones de alcance.} La ausencia inicial de
autenticación fue la que generó más incomodidad real durante el desarrollo: saber que,
mientras no se implementara, cualquiera con acceso a la API podía consultar o
modificar los datos de cualquier cliente sin restricción, se sentía como una brecha
real y no solo teórica -- razón por la cual se decidió cerrarla al menos parcialmente
antes de la entrega final, en lugar de dejarla únicamente documentada como limitación.

Una duda similar surgió con la exclusión de reportes de venta y de la gestión de roles
administrativos (sección~\ref{ch:arquitectura}): un Jefe de Sucursal o un Administrador
del Sistema reales necesitarían esas funciones para operar el negocio día a día, y su
ausencia no es solo una limitación técnica, sino una función de negocio real que el
caso pide. A diferencia de la autenticación, esta duda llevó a revisar el alcance
antes de la entrega final: los reportes de ventas por sucursal y por autor se
implementaron (sección~\ref{ch:crud}), por ser acotados y no requerir cambios de
arquitectura. La gestión de usuarios y permisos por rol, en cambio, se decidió
dejar fuera deliberadamente -- a diferencia de un reporte, que es una vista de solo
lectura sobre datos que ya existen, un sistema de roles reales implica un modelo de
autorización completo y una decisión arquitectónica (¿es un nuevo microservicio? ¿se
integra a uno existente?) que merece más tiempo de análisis del que quedaba
disponible, y apurarla habría sido peor que documentarla honestamente como
limitación.
